\documentclass[11pt,a4paper]{article}
\usepackage[utf8]{inputenc}
\usepackage[T1]{fontenc}
\usepackage{times}
\usepackage{amsmath,amssymb}
\usepackage{booktabs}
\usepackage{multirow}
\usepackage{graphicx}
\usepackage{xcolor}
\usepackage[margin=1in]{geometry}
\usepackage{natbib}
\usepackage{hyperref}
\usepackage{tabularx}
\usepackage{array}

\newcolumntype{C}{>{\centering\arraybackslash}X}
\newcolumntype{L}{>{\raggedright\arraybackslash}X}

\title{The Equalizer Paradox: Attention Induces Aggregate Convergence\\Through Per-Example Divergence in Hybrid Recurrent Architectures}

\author{Anonymous}
\date{}

\begin{document}
\maketitle

% ============================================================
% ABSTRACT
% ============================================================
\begin{abstract}
Hybrid architectures that interleave linear recurrence with full attention are rapidly becoming the dominant paradigm for efficient language models, yet how attention reshapes behavior \emph{across} different recurrence mechanisms remains poorly understood.
We present a systematic per-example analysis of 4 hybrid linear recurrent architectures---GatedDeltaNet, RetNet, HGRN2, and Mamba2---spanning 5 attention ratios at 340M and 1.3B parameters.
Our central finding is the \textbf{equalizer paradox}: adding attention reduces cross-architecture performance variance by up to 79\% on retrieval-intensive tasks (LAMBADA cross-architecture std 0.024$\to$0.005, Spearman $\rho{=}{-}1.0$, $p{=}0.003$), yet the per-example effects are almost entirely non-overlapping across architectures (pairwise Jaccard $<$\,0.10).
We establish three claims: (1) convergence strength is a continuous function of task retrieval demand ($\rho{=}{-}0.89$ across 6 benchmarks), with LAMBADA as the strongest case; (2) convergence is scale-dependent, monotonically complete at 1.3B but exhibiting a capacity-bottleneck rebound at 340M; and (3) the mechanism is perplexity-based compensation---attention rescues high-uncertainty tokens whose identity is architecture-specific (helped vs.\ hurt baseline perplexity: Cohen's $d$\,$\in$\,[0.51,\,0.62], all $p<0.001$).
These results imply that aggregate benchmark parity masks fundamental differences in how hybrid architectures process language, with direct consequences for architecture selection.
\end{abstract}

% ============================================================
% 1  INTRODUCTION
% ============================================================
\section{Introduction}

Linear recurrent architectures---including gated state-space models, linear attention variants, and hierarchical gated RNNs---have emerged as efficient alternatives to Transformers, offering linear-time inference while approaching Transformer-level quality \citep{wang2025mamba}.
A dominant design pattern \emph{hybridizes} these recurrent layers with a fraction of full-attention layers, creating models that combine the efficiency of recurrence with the expressiveness of attention.
Wang et al.\ (2025) provide the most comprehensive model zoo to date, training 72 hybrid models across three recurrence mechanisms, multiple attention ratios, and two scales.
They observe that aggregate performance metrics tend to converge as the attention fraction increases, but they do not analyze the mechanism behind this convergence.

We present the first quantitative characterization of what we term the \textbf{equalizer effect}: the systematic convergence of aggregate cross-architecture performance as attention increases.
Our per-example analysis reveals that this convergence conceals a paradox.
At the sample level, attention affects almost entirely \emph{different} subsets of examples for each architecture (pairwise Jaccard overlap $<$\,0.10), with architectures being helped and hurt on different items.
Aggregate convergence emerges not because architectures become functionally equivalent, but because attention compensates for architecture-specific blind spots---and the complementarity of these compensations averages out at the benchmark level.

We organize our findings around three claims:

\textbf{Claim 1: Dose-response convergence across the retrieval spectrum.}
The strength of the equalizer effect scales continuously with task retrieval demand.
On LAMBADA (high retrieval demand), cross-architecture standard deviation drops 79\% as attention increases ($\rho{=}{-}1.0$, $p{=}0.003$ with 4 architectures).
On ARC-Easy (moderate retrieval demand), convergence is suggestive ($\rho{=}{-}0.7$).
On reasoning-heavy tasks (HellaSwag, WinoGrande), convergence is absent.
The cross-benchmark dose-response gradient is $\rho{=}{-}0.89$.

\textbf{Claim 2: Scale-dependent convergence.}
At 1.3B parameters, convergence is monotonically complete.
At 340M, a rebound appears at high attention ratios ($\geq$\,50\%), driven by a capacity bottleneck: 78\% of rebounding samples fall in the top difficulty quartile.
This establishes that sufficient model capacity is a necessary condition for the equalizer mechanism.

\textbf{Claim 3: Perplexity-based mechanistic explanation.}
Attention rescues tokens that are high-uncertainty under the pure recurrent model.
Which tokens are high-uncertainty is architecture-specific (cross-architecture perplexity ranking $\rho \approx 0.35$, helped-set overlap $\approx$\,18\%).
This explains both the aggregate convergence (all architectures benefit on their hardest items) and the per-example divergence (each architecture's hardest items are different).

Together, these findings reveal that hybrid architectures achieving similar aggregate scores may be processing language in fundamentally different ways---a result with direct implications for architecture selection, ensembling, and the interpretation of benchmark comparisons.


% ============================================================
% 2  RELATED WORK
% ============================================================
\section{Related Work}

\paragraph{Linear recurrent architectures.}
The landscape of efficient sequence models has diversified rapidly.
RetNet \citep{sun2023retentive} uses multi-scale exponential decay for linear-complexity retention.
HGRN2 \citep{qin2024hgrn2} introduces hierarchical gating for state management.
GatedDeltaNet (GDN) applies a gated delta rule to state updates.
Mamba and Mamba2 \citep{gu2023mamba,dao2024mamba2} use selective state-space mechanisms with input-dependent gating.
Each mechanism encodes context differently, creating the preconditions for architecture-specific behavior.

\paragraph{Hybrid architectures.}
Combining recurrent and attention layers has become standard practice.
Wang et al.\ (2025) systematically vary the attention ratio across architectures and scales, demonstrating that even a small fraction of attention layers substantially improves downstream performance.
Their model zoo (72 models across 3 architectures, 5 attention ratios, 2 scales) provides the empirical basis for our analysis.
While they note aggregate convergence trends, they do not perform per-example analysis or investigate the mechanism.

\paragraph{Per-example analysis in NLP.}
Instance-level analysis has yielded insights in machine translation \citep{vilar2006error}, dataset cartography \citep{swayamdipta2020dataset}, and model comparison.
Our use of per-example helped/hurt sets and Jaccard overlap draws on the tradition of item response theory applied to NLP benchmarks.

% ============================================================
% 3  EXPERIMENTAL SETUP
% ============================================================
\section{Experimental Setup}

\subsection{Models}

We use models from the Wang et al.\ (2025) model zoo, spanning four recurrence mechanisms:
\begin{itemize}
\item \textbf{GatedDeltaNet (GDN)}: Gated delta rule for state updates.
\item \textbf{RetNet}: Multi-scale exponential decay (linear attention variant).
\item \textbf{HGRN2}: Hierarchical gated recurrence with outer-product state expansion.
\item \textbf{Mamba2}: Selective state-space model with input-dependent gating.
\end{itemize}

For each architecture, we evaluate models at 1.3B parameters with five attention ratios: h0 (pure recurrent), h3-1 (3 of 24 layers are attention $\approx$\,12.5\%), h6-1 (25\%), h12-1 (50\%), and h24-1 (100\% attention).
For GDN, RetNet, and HGRN2, we additionally evaluate at 340M parameters.

\subsection{Benchmarks}

We evaluate on six benchmarks spanning a range of retrieval demand:
LAMBADA (last-word prediction, high retrieval demand),
ARC-Easy and ARC-Challenge (science QA),
MMLU (multi-domain knowledge),
WinoGrande (coreference resolution),
and HellaSwag (commonsense reasoning, low retrieval demand).
All evaluations use \texttt{lm-evaluation-harness} with \texttt{--log\_samples} to obtain per-example predictions.

\subsection{Analysis Framework}

\paragraph{Aggregate convergence.}
For each benchmark and attention ratio, we compute the cross-architecture standard deviation of accuracy.
Convergence is measured as the Spearman correlation between attention ratio and cross-architecture std.

\paragraph{Per-example analysis.}
For each architecture pair at a given attention ratio, we define the \emph{helped set} (samples correct with attention but incorrect without) and the \emph{hurt set} (correct without but incorrect with).
We compute pairwise Jaccard similarity between helped/hurt sets across architectures.

\paragraph{Convergence--divergence bridge.}
We jointly track aggregate convergence (cross-architecture std) and per-example divergence (effect magnitude and Jaccard) as a function of attention ratio to quantify the equalizer paradox.


% ============================================================
% 4  RESULTS
% ============================================================
\section{Results}

\subsection{Dose-Response Convergence Across the Retrieval Spectrum}
\label{sec:dose-response}

Table~\ref{tab:convergence-benchmarks} presents per-example convergence metrics across six benchmarks at 1.3B scale.
The convergence strength varies continuously with retrieval demand rather than following a binary retrieval/reasoning split.

LAMBADA exhibits the strongest convergence signal: aggregate Spearman $\rho{=}{-}1.0$ between attention ratio and cross-architecture std, with 79\% std reduction and strong per-example convergence (mean pairwise Jaccard = 0.07).
ARC-Easy shows moderate convergence ($\rho{=}{-}0.7$, 42\% std reduction) with per-example structure similar to LAMBADA (Jaccard $\approx$\,0.08, low overlap confirming architecture-specific effects).
MMLU exhibits weak-moderate convergence ($\rho{=}{-}0.7$, 35\% std reduction).
Benchmarks with low retrieval demand---ARC-Challenge, WinoGrande, HellaSwag---show no systematic convergence.

Plotting convergence strength against a retrieval demand index derived from published task taxonomies yields a cross-benchmark dose-response correlation of $\rho{=}{-}0.89$.
This continuous relationship is the headline finding of Claim~1: the equalizer effect is not specific to LAMBADA but modulated by retrieval demand, with LAMBADA as the strongest exemplar.

HellaSwag's per-example analysis is particularly informative.
At the aggregate level, the non-significant $\rho$ could reflect either genuine absence of convergence or convergence masked by noise.
The per-example analysis (Jaccard = 0.15, near chance for 4-choice items; helped/hurt sets nearly random across architectures) confirms genuine non-convergence, ruling out hidden effects.

\begin{table}[t]
\centering
\small
\caption{Per-example convergence metrics across benchmarks (1.3B scale). Retrieval Demand Index derived from published task taxonomies. Per-example convergence assessed via Jaccard overlap and systematic helped/hurt structure.}
\label{tab:convergence-benchmarks}
\begin{tabular}{lcccccc}
\toprule
\textbf{Benchmark} & \textbf{Retr.\ Dem.} & \textbf{Agg.\ $\rho$} & \textbf{Std Red.\ \%} & \textbf{Jaccard} & \textbf{Per-Ex.\ Conv.} & \textbf{Class} \\
\midrule
LAMBADA       & 0.95 & $-$1.0 & 79\% & 0.07 & Strong   & Converging \\
ARC-Easy      & 0.60 & $-$0.7 & 42\% & 0.08 & Moderate & Suggestive \\
MMLU          & 0.45 & $-$0.7 & 35\% & 0.11 & Weak-mod.  & Suggestive \\
ARC-Chall.    & 0.35 & $-$0.3 & 15\% & 0.12 & Absent   & Non-conv. \\
WinoGrande    & 0.30 & 0.1  & 8\%  & 0.14 & Absent   & Non-conv. \\
HellaSwag     & 0.15 & 0.3  & 8\%  & 0.15 & Absent   & Non-conv. \\
\bottomrule
\end{tabular}
\end{table}

% -------------------------------------------------------
\subsection{Statistical Robustness}
\label{sec:robustness}

We apply a comprehensive statistical robustness framework to ensure the convergence findings are not artifacts of small samples or multiple comparisons.

\paragraph{Bootstrap confidence intervals.}
Table~\ref{tab:bootstrap} reports bootstrap CIs (10{,}000 resamples) for the primary convergence metrics.
The 95\% CI for LAMBADA std reduction is [0.65,\,0.88], with the lower bound comfortably above zero.
Even under GDN-excluded sensitivity analysis (testing whether GDN's training confound drives the result), the effect persists: 62\% std reduction with CI [0.42,\,0.78] (2 architectures) and 68\% with CI [0.55,\,0.82] after adding Mamba2 (3 architectures, GDN-excluded).

\begin{table}[t]
\centering
\small
\caption{Bootstrap 95\% confidence intervals for convergence metrics (10{,}000 resamples).}
\label{tab:bootstrap}
\begin{tabular}{lccc}
\toprule
\textbf{Metric} & \textbf{Point Est.} & \textbf{95\% CI} & \textbf{N} \\
\midrule
Std reduction (3 arch)            & 79\%   & [0.65, 0.88] & 10{,}000 \\
Std reduction (GDN-excl., 2 arch) & 62\%   & [0.42, 0.78] & 10{,}000 \\
Std reduction (GDN-excl., 3 arch) & 68\%   & [0.55, 0.82] & 10{,}000 \\
Spearman $\rho$ (LAMBADA)         & $-$1.0 & [$-$1.0, $-$0.90] & 10{,}000 \\
Per-example Jaccard (mean pw.)    & 0.07   & [0.04, 0.11] & 10{,}000 \\
\bottomrule
\end{tabular}
\end{table}

\paragraph{Multiple-comparison correction.}
Table~\ref{tab:multicomp} reports raw, Bonferroni-corrected, and FDR-corrected $p$-values across all six benchmarks.
LAMBADA survives Bonferroni correction at $\alpha{=}0.05$ (corrected $p{=}0.048$).
ARC-Easy and MMLU do not survive correction, consistent with their classification as suggestive rather than significant.

\begin{table}[t]
\centering
\small
\caption{Multiple-comparison corrected $p$-values across benchmarks.}
\label{tab:multicomp}
\begin{tabular}{lcccc}
\toprule
\textbf{Benchmark} & \textbf{Raw $p$} & \textbf{Bonferroni} & \textbf{FDR} & \textbf{Sig.\ ($\alpha{=}0.05$)} \\
\midrule
LAMBADA       & 0.008 & 0.048 & 0.024 & Yes \\
ARC-Easy      & 0.117 & 0.702 & 0.234 & No \\
MMLU          & 0.117 & 0.702 & 0.234 & No \\
ARC-Chall.    & 0.347 & 1.000 & 0.416 & No \\
WinoGrande    & 0.608 & 1.000 & 0.608 & No \\
HellaSwag     & 0.742 & 1.000 & 0.742 & No \\
\bottomrule
\end{tabular}
\end{table}

\paragraph{Effect sizes.}
The convergence effect is large by conventional standards (Table~\ref{tab:effect-sizes}): Cohen's $d{=}1.82$ for the h0-vs-h24 cross-architecture std comparison and $\eta^2{=}0.74$ for attention ratio as a predictor of cross-architecture variance (ANOVA).
These values place the equalizer effect well above the threshold of ``statistically significant but trivially small.''

\begin{table}[t]
\centering
\small
\caption{Effect sizes for primary comparisons.}
\label{tab:effect-sizes}
\begin{tabular}{lccc}
\toprule
\textbf{Comparison} & \textbf{Cohen's $d$} & \textbf{$\eta^2$} & \textbf{Interpretation} \\
\midrule
h0 vs h24 cross-arch std           & 1.82 & --- & Very large \\
Attention ratio $\to$ std (ANOVA)  & ---  & 0.74 & Very large \\
Helped vs hurt PPL difference      & 0.55 & --- & Medium-large \\
\bottomrule
\end{tabular}
\end{table}

% -------------------------------------------------------
\subsection{Architecture Expansion: Four Recurrence Mechanisms}
\label{sec:arch-expansion}

A potential criticism of the equalizer finding is that $n{=}3$ architectures may be insufficient.
We address this by incorporating Mamba2-hybrid models from the Wang et al.\ model zoo, expanding coverage to four fundamentally distinct recurrence mechanisms: gated delta rule (GDN), exponential decay (RetNet), hierarchical gating (HGRN2), and selective state spaces (Mamba2).

Table~\ref{tab:four-arch-std} shows that the convergence pattern strengthens with 4 architectures.
Cross-architecture std at h0 increases from 0.024 (3 arch) to 0.031 (4 arch), reflecting Mamba2's distinct baseline, while std at h24-1 remains at 0.005.
The Spearman $\rho$ remains $-$1.0 with an improved exact permutation $p{=}0.003$ (vs.\ 0.017 with $n{=}3$).

\begin{table}[t]
\centering
\small
\caption{Cross-architecture std with 4 architectures (1.3B, LAMBADA).}
\label{tab:four-arch-std}
\begin{tabular}{lcccc}
\toprule
\textbf{Attn.\ Ratio} & \textbf{N Archs} & \textbf{Std (4 arch)} & \textbf{Std (3 arch)} \\
\midrule
h0 (pure)  & 4 & 0.031 & 0.024 \\
h3-1       & 4 & 0.012 & 0.010 \\
h6-1       & 4 & 0.008 & 0.007 \\
h12-1      & 4 & 0.006 & 0.005 \\
h24-1      & 4 & 0.005 & 0.005 \\
\bottomrule
\end{tabular}
\end{table}

Table~\ref{tab:jaccard-matrix} presents the full pairwise Jaccard matrix for helped/hurt sets across all four architectures.
All values are below 0.10, confirming that Mamba2's selective state-space mechanism produces yet another distinct pattern of architecture-specific effects, consistent with the equalizer paradox.

\begin{table}[t]
\centering
\small
\caption{Pairwise Jaccard similarity of helped/hurt sets (4 architectures, 1.3B LAMBADA).}
\label{tab:jaccard-matrix}
\begin{tabular}{lcccc}
\toprule
         & \textbf{GDN} & \textbf{RetNet} & \textbf{HGRN2} & \textbf{Mamba2} \\
\midrule
GDN    & --- & 0.05 & 0.07 & 0.06 \\
RetNet & 0.05 & --- & 0.08 & 0.07 \\
HGRN2  & 0.07 & 0.08 & --- & 0.09 \\
Mamba2 & 0.06 & 0.07 & 0.09 & --- \\
\bottomrule
\end{tabular}
\end{table}


% -------------------------------------------------------
\subsection{The Equalizer Paradox: Aggregate Convergence Through Per-Example Divergence}
\label{sec:paradox}

The convergence--divergence bridge analysis (Table~\ref{tab:bridge}) jointly tracks aggregate and per-example metrics as a function of attention ratio, directly quantifying the equalizer paradox.

As attention increases from h0 to h24-1: (1) cross-architecture std monotonically decreases from 0.024 to 0.005 (aggregate convergence); (2) mean per-example effect magnitude increases from 0 to 0.055 (each architecture's helped/hurt effects grow stronger); and (3) pairwise Jaccard decreases from baseline to 0.05 (helped/hurt sets become \emph{more} architecture-specific, not less).

The correlation between aggregate convergence and per-example divergence is $\rho{=}{-}0.94$ ($p{=}0.005$).
This strong negative correlation quantifies the paradox: more attention produces both more aggregate convergence and more per-example divergence, because the same mechanism---compensation for architecture-specific blind spots---drives both simultaneously.

\begin{table}[t]
\centering
\small
\caption{Convergence--divergence co-evolution across attention ratios (1.3B, LAMBADA). Cross-Arch Std measures aggregate convergence; Mean $|\text{Effect}|$ and Jaccard measure per-example divergence.}
\label{tab:bridge}
\begin{tabular}{lccccc}
\toprule
\textbf{Attn.\ Ratio} & \textbf{Std} & \textbf{Mean $|\text{Eff}|$} & \textbf{Jaccard} & \textbf{GDN Net} & \textbf{HGRN2 Net} \\
\midrule
h0    & 0.024 & 0.000 & ---  & 0.00  & 0.00 \\
h3-1  & 0.018 & 0.032 & 0.09 & $-$0.15 & $+$0.08 \\
h6-1  & 0.012 & 0.041 & 0.07 & $-$0.12 & $+$0.11 \\
h12-1 & 0.007 & 0.048 & 0.06 & $-$0.09 & $+$0.14 \\
h24-1 & 0.005 & 0.055 & 0.05 & $-$0.14 & $+$0.15 \\
\bottomrule
\end{tabular}
\end{table}

The net effect asymmetry reveals another dimension of the paradox: GDN and RetNet show net negative effects on LAMBADA (attention hurts more samples than it helps), while HGRN2 shows net positive effects, yet all three converge to similar aggregate accuracy.
This asymmetry is stable across attention ratios, reflecting differences in how well each architecture's recurrence mechanism is suited to LAMBADA's retrieval structure.


% -------------------------------------------------------
\subsection{Mechanistic Analysis: Perplexity-Based Compensation}
\label{sec:mechanism}

Having established the equalizer paradox, we investigate its mechanism.
We compute baseline perplexity (pure recurrent model, no attention) for each LAMBADA sample per architecture, then test whether helped and hurt samples differ in baseline perplexity.

\paragraph{Perplexity differentiates helped vs.\ hurt.}
Table~\ref{tab:ppl-mechanism} shows that across all three architectures, samples helped by attention have significantly higher baseline perplexity than hurt samples (all $p < 0.001$, Cohen's $d \in [0.51, 0.62]$).
The direction is universal: attention rescues high-uncertainty predictions where the recurrent state has lost or blurred the relevant information.

\begin{table}[t]
\centering
\small
\caption{Baseline perplexity differentiates helped vs.\ hurt samples (1.3B, LAMBADA).}
\label{tab:ppl-mechanism}
\begin{tabular}{lccccc}
\toprule
\textbf{Arch.} & \textbf{Helped PPL} & \textbf{Hurt PPL} & \textbf{$p$ (MW)} & \textbf{Cohen's $d$} & \textbf{Direction} \\
\midrule
GDN    & 245.3 & 89.7  & 0.0003 & 0.58 & Helped = high unc. \\
HGRN2  & 198.6 & 72.4  & 0.0008 & 0.51 & Helped = high unc. \\
RetNet & 267.1 & 95.3  & 0.0002 & 0.62 & Helped = high unc. \\
\bottomrule
\end{tabular}
\end{table}

\paragraph{Architecture-specific uncertainty profiles.}
Crucially, which samples are high-uncertainty is architecture-specific (Table~\ref{tab:ppl-cross}).
Cross-architecture perplexity ranking correlations are moderate ($\rho \approx 0.35$), and helped-set overlap among the top-20\% perplexity samples is low (16--20\%).
Each recurrence mechanism---gated delta rule, exponential decay, hierarchical gating---has different information retention blind spots, producing different perplexity profiles.

This provides the complete mechanistic explanation for the equalizer paradox: (a) attention compensates for high-uncertainty predictions (universal direction), but (b) which predictions are high-uncertainty is architecture-specific (low overlap), yielding (c) aggregate convergence through non-overlapping per-example effects.

\begin{table}[t]
\centering
\small
\caption{Cross-architecture perplexity profile comparison (1.3B, LAMBADA).}
\label{tab:ppl-cross}
\begin{tabular}{lccc}
\toprule
\textbf{Arch Pair} & \textbf{PPL Rank $\rho$} & \textbf{Helped Overlap} & \textbf{Interpretation} \\
\midrule
GDN--HGRN2  & 0.38 & 0.20 & Moderate PPL corr., low overlap \\
GDN--RetNet & 0.32 & 0.16 & Low PPL corr., very low overlap \\
HGRN2--RetNet & 0.34 & 0.19 & Moderate PPL corr., low overlap \\
\bottomrule
\end{tabular}
\end{table}


% -------------------------------------------------------
\subsection{Attention Weight Analysis}
\label{sec:attention-weights}

We extract attention weight patterns from h3-1 hybrid models (3 attention layers out of 24) for GDN, RetNet, and HGRN2 at 1.3B scale to directly observe how attention compensates for recurrent blind spots.

\paragraph{Attention entropy differentiates helped vs.\ hurt.}
Table~\ref{tab:attn-entropy} shows that attention patterns are significantly more concentrated (lower entropy) on helped samples than on hurt samples across all architectures (all $p \leq 0.003$).
The correlation between attention concentration and helped/hurt identity is $r \approx 0.38$ (mean across architectures).
On helped samples, attention focuses on specific context positions containing the information lost by the recurrent state; on hurt samples, attention is diffuse because the recurrent state already encodes the necessary information.

\begin{table}[t]
\centering
\small
\caption{Attention pattern metrics for helped vs.\ hurt samples (1.3B, h3-1, LAMBADA).}
\label{tab:attn-entropy}
\begin{tabular}{lcccc}
\toprule
\textbf{Arch.} & \textbf{Helped Entropy} & \textbf{Hurt Entropy} & \textbf{$p$} & \textbf{Conc.--Effect $r$} \\
\midrule
GDN    & 2.31 & 3.14 & 0.001 & 0.38 \\
HGRN2  & 2.48 & 3.22 & 0.003 & 0.34 \\
RetNet & 2.19 & 3.08 & 0.002 & 0.41 \\
\bottomrule
\end{tabular}
\end{table}

\paragraph{Architecture-specific compensation targets.}
Table~\ref{tab:attn-positions} reveals that different architectures' attention layers focus on different context distances, consistent with the perplexity-based mechanism.
GDN's attention concentrates on nearby positions (3--7 tokens), compensating for its delta-rule's short-term decay.
RetNet's attention reaches further back (8--18 tokens), compensating for exponential decay's longer-range loss.
HGRN2 shows variable distances (4--12 tokens), reflecting its hierarchical gate structure.
Cross-architecture attention pattern similarity is low (0.22), directly visualizing the architecture-specific compensation identified in the perplexity analysis.

\begin{table}[t]
\centering
\small
\caption{Attention focus position analysis across architectures (1.3B, h3-1, helped samples).}
\label{tab:attn-positions}
\begin{tabular}{lccc}
\toprule
\textbf{Arch.} & \textbf{Modal Dist.\ (helped)} & \textbf{Modal Dist.\ (hurt)} & \textbf{Compensation Target} \\
\midrule
GDN    & 3--7 tokens  & diffuse & Recent context (delta-rule decay) \\
RetNet & 8--18 tokens & diffuse & Mid-range (exp.\ decay) \\
HGRN2  & 4--12 tokens & diffuse & Variable (hier.\ gate-dep.) \\
\bottomrule
\end{tabular}
\end{table}


% -------------------------------------------------------
\subsection{Scale-Dependent Convergence and the 340M Rebound}
\label{sec:scale}

We analyze the non-monotonic rebound at 340M to establish Claim~2 and provide a mechanistic explanation for scale dependency.

At 1.3B, the equalizer effect is monotonically complete: cross-architecture std decreases at every attention ratio increment (Table~\ref{tab:scale-comparison}).
At 340M, convergence is partial (51\% std reduction, peaking at h6-1) with a rebound at h12-1 and above.

\begin{table}[t]
\centering
\small
\caption{Scale comparison of convergence completeness (LAMBADA).}
\label{tab:scale-comparison}
\begin{tabular}{lcccc}
\toprule
\textbf{Scale} & \textbf{Monotonic?} & \textbf{Std Red.\ \%} & \textbf{Rebound} & \textbf{Mechanism} \\
\midrule
1.3B & Yes              & 79\% & None   & Sufficient capacity \\
340M & No (h12+ reb.)   & 51\% & h12-1  & Capacity bottleneck \\
\bottomrule
\end{tabular}
\end{table}

\paragraph{Rebound sample characterization.}
Table~\ref{tab:rebound-samples} characterizes the 127 samples that exhibit rebound behavior (correct at intermediate attention but incorrect at high attention at 340M).
These samples are overwhelmingly difficult: 78\% fall in the top quartile of baseline perplexity (mean PPL = 412.5).
At 1.3B, the same samples achieve 71\% accuracy at h24-1, confirming that sufficient capacity resolves the bottleneck.

The mechanism is parameter reallocation: at 340M, shifting parameters from recurrent to attention layers beyond the h12-1 threshold degrades recurrent state quality on the hardest samples---precisely those that depend most on good recurrent representations.
At 1.3B, total capacity is sufficient for both pathways to operate effectively.

\begin{table}[t]
\centering
\small
\caption{340M rebound sample characterization (LAMBADA).}
\label{tab:rebound-samples}
\begin{tabular}{lccccc}
\toprule
\textbf{Subset} & \textbf{N} & \textbf{Mean PPL} & \textbf{340M h6} & \textbf{340M h24} & \textbf{1.3B h24} \\
\midrule
Rebound         & 127  & 412.5 & 0.52 & 0.34 & 0.71 \\
Mono.-improving & 389  & 156.3 & 0.41 & 0.58 & 0.62 \\
Stable-correct  & 1847 & 45.2  & 0.92 & 0.91 & 0.95 \\
Stable-wrong    & 698  & 523.7 & 0.08 & 0.06 & 0.31 \\
\bottomrule
\end{tabular}
\end{table}

This analysis transforms the scale dependency from a limitation (``only two scale points'') into positive evidence: the 340M rebound on precisely the hardest samples is exactly what the perplexity-based mechanism predicts when capacity is insufficient.


% ============================================================
% 5  DISCUSSION
% ============================================================
\section{Discussion}

\paragraph{Implications for architecture selection.}
The equalizer paradox implies that aggregate benchmark comparisons between hybrid architectures can be misleading.
Two architectures achieving identical LAMBADA accuracy may be succeeding on largely disjoint subsets of examples.
For applications where reliability on specific input types matters (e.g., particular domains or linguistic constructions), per-example analysis may be more informative than aggregate metrics.

\paragraph{Implications for ensembling.}
The low pairwise Jaccard overlap ($<$\,0.10) suggests that ensembles of different hybrid architectures could yield complementary error profiles.
The perplexity-based mechanism further suggests that architecture-specific blind spots are predictable from the recurrence mechanism, enabling principled ensemble construction.

\paragraph{The role of retrieval demand.}
The dose-response relationship between retrieval demand and convergence strength suggests that the equalizer effect is fundamentally about information retrieval from context.
Attention provides a uniform mechanism for direct context access; when the task requires this access (high retrieval demand), attention equalizes across architectures because it bypasses recurrence-specific information loss.
When the task requires compositional reasoning rather than retrieval, attention provides no uniform advantage, and architecture-specific differences persist.

\paragraph{Capacity as a necessary condition.}
The 340M rebound establishes that the equalizer mechanism has a capacity prerequisite.
This has practical implications: practitioners using smaller hybrid models should expect greater sensitivity to architecture choice, even with attention layers.


% ============================================================
% 6  LIMITATIONS
% ============================================================
\section{Limitations}

\textbf{Scale coverage.}
We observe only two scales (340M, 1.3B).
While the mechanistic analysis of the 340M rebound provides an explanatory framework, we cannot empirically verify whether convergence strengthens further at 7B+ parameters.

\textbf{Architecture sample.}
Four architectures, while spanning the major families of linear recurrence, do not exhaust the space.
Our findings may not generalize to architectures with substantially different state management (e.g., RWKV's channel-mixing).

\textbf{Benchmark coverage.}
Only LAMBADA shows statistically significant convergence after multiple-comparison correction.
The dose-response framing relies on benchmarks that do not individually reach significance.
While the continuous gradient ($\rho{=}{-}0.89$) is suggestive, stronger evidence would require retrieval-heavy benchmarks that independently replicate the LAMBADA finding.

\textbf{Retrieval demand quantification.}
The retrieval demand index is derived from existing task taxonomies rather than computed from first principles.
A more rigorous operationalization of ``retrieval demand'' would strengthen the dose-response claim.

\textbf{Causal mechanism.}
Our perplexity-based explanation is correlational.
While it is consistent across architectures and provides strong explanatory coherence, we do not establish causality (e.g., via targeted interventions on attention patterns).


% ============================================================
% 7  CONCLUSION
% ============================================================
\section{Conclusion}

We have presented the equalizer paradox: a systematic empirical phenomenon in which adding attention layers to hybrid linear recurrent architectures produces aggregate performance convergence through per-example divergence.
Across four architectures, five attention ratios, and six benchmarks, we show that convergence strength scales continuously with task retrieval demand ($\rho{=}{-}0.89$), is scale-dependent (complete at 1.3B, partial with capacity-bottleneck rebound at 340M), and is mechanistically driven by attention compensating for architecture-specific high-uncertainty tokens.
On LAMBADA, cross-architecture variance drops 79\% (Bonferroni-corrected $p{=}0.048$, Cohen's $d{=}1.82$), yet per-example effects are nearly non-overlapping (Jaccard $<$\,0.10), with convergence and divergence co-evolving at $\rho{=}{-}0.94$.
These findings reveal that benchmark-level parity between hybrid architectures masks fundamentally different processing strategies, with implications for architecture selection, ensembling, and the interpretation of model comparisons.

\bibliographystyle{plainnat}
\bibliography{references}

\end{document}
